(1) 纯物质液体与蒸气在压强 \(P\)、温度 \(T\) 下共存，两相摩尔体积为 \(v_l,v_g\)，摩尔熵为 \(s_l,s_g\)。从化学势两相相等推导共存线的 \(dP/dT\)。
(2) 以摩尔汽化潜热 \(L=T(s_g-s_l)\) 表示斜率，说明 \(L>0,v_g>v_l\) 时共存线的方向。
(3) 假设蒸气可视为理想气体、液体体积可忽略、潜热 \(L\) 近似常数，积分得到两温度间饱和蒸气压之比。
(4) 用绝对温度与压强的极限说明为什么高温临界区域不能无限延拓该公式，并检查指数的无量纲性。
